% GENERATED FILE. Edit canonical lessons, then run npm run generate:books.
% canonical-source-hash: fnv1a64:012064b689551da1
% canonical-lessons: KA-C137-meju, KA-C137-kitaki, KA-C137-kannadi, KA-C137-tivi, KA-C137-kampyutar

\chapter{Table, Window, Mirror, Television, Computer}
\label{ch:ka-table-window-mirror-television-computer}
\hlchaptermodality{137}

\begin{chapteropening}
\textbf{By the end of this chapter:} \emph{I can say a table, a window, a mirror, a television and a computer.}

Complete the last lesson of chapter 137: I can say a table, a window, a mirror, a television and a computer.
\end{chapteropening}

\section[\texorpdfstring{\kn{ಮೇಜು} (mēju)}{ಮೇಜು (mēju)}]{\kn{ಮೇಜು} (mēju) — a table}

\label{lesson:KA-C137-meju}



\begin{quote}
Before the new one: say the Kannada for the waist, then the Kannada for an eyebrow.
\end{quote}

\subsection*{You'll want to know: \kn{ಮೇಜು}}
\textbf{\kn{ಮೇಜು}} — \emph{mēju} — \textquotedblleft{}a table\textquotedblright{}.

\begin{grammarlens}[title={What you've built}]
One more word for this chapter.
\end{grammarlens}

\subsection*{Your turn}
\begin{itemize}
\raggedright
  \item \emph{Say it:} \emph{mēju}
  \item \emph{Say it:} \emph{mēju}, once more
  \item \emph{Say it:} say \emph{mēju}
  \item \emph{Recall:} say the Kannada for the waist, then the Kannada for an eyebrow, then say \emph{mēju} again
\end{itemize}

\begin{tcolorbox}[breakable,colback=teal!4,colframe=teal!35!black,title={Before you move on}]
What does \kn{ಮೇಜು} mean? (\textquotedblleft{}A table\textquotedblright{}.) Say it once more.
\end{tcolorbox}
\section[\texorpdfstring{\kn{ಕಿಟಕಿ} (kiṭaki)}{ಕಿಟಕಿ (kiṭaki)}]{\kn{ಕಿಟಕಿ} (kiṭaki) — a window}

\label{lesson:KA-C137-kitaki}



\begin{quote}
Before the new one: say the Kannada for an eyebrow, then the Kannada for a table.

\emph{Read it:} \textbf{\kn{ಮರ} \kn{ಅಲ್ಲಿ} \kn{ಇದೆ}. \kn{ನದಿ} \kn{ಅಲ್ಲಿ} \kn{ಇದೆ}.} aloud as one run
\end{quote}

\subsection*{You'll want to know: \kn{ಕಿಟಕಿ}}
\textbf{\kn{ಕಿಟಕಿ}} — \emph{kiṭaki} — \textquotedblleft{}a window\textquotedblright{}.

\begin{grammarlens}[title={What you've built}]
One more word for this chapter.
\end{grammarlens}

\subsection*{Your turn}
\begin{itemize}
\raggedright
  \item \emph{Say it:} \emph{kiṭaki}
  \item \emph{Say it:} \emph{kiṭaki}, once more
  \item \emph{Say it:} say \emph{kiṭaki}
  \item \emph{Recall:} say the Kannada for an eyebrow, then the Kannada for a table, then say \emph{kiṭaki} again
\end{itemize}

\begin{tcolorbox}[breakable,colback=teal!4,colframe=teal!35!black,title={Before you move on}]
What does \kn{ಕಿಟಕಿ} mean? (\textquotedblleft{}A window\textquotedblright{}.) Say it once more.
\end{tcolorbox}
\section[\texorpdfstring{\kn{ಕನ್ನಡಿ} (kannaḍi)}{ಕನ್ನಡಿ (kannaḍi)}]{\kn{ಕನ್ನಡಿ} (kannaḍi) — a mirror}

\label{lesson:KA-C137-kannadi}



\begin{quote}
Before the new one: say the Kannada for a table, then the Kannada for a window.
\end{quote}

\subsection*{You'll want to know: \kn{ಕನ್ನಡಿ}}
\textbf{\kn{ಕನ್ನಡಿ}} — \emph{kannaḍi} — \textquotedblleft{}a mirror\textquotedblright{}.

\begin{grammarlens}[title={What you've built}]
One more word for this chapter.
\end{grammarlens}

\subsection*{Your turn}
\begin{itemize}
\raggedright
  \item \emph{Say it:} \emph{kannaḍi}
  \item \emph{Say it:} \emph{kannaḍi}, once more
  \item \emph{Say it:} say \emph{kannaḍi}
  \item \emph{Recall:} say the Kannada for a table, then the Kannada for a window, then say \emph{kannaḍi} again
\end{itemize}

\begin{tcolorbox}[breakable,colback=teal!4,colframe=teal!35!black,title={Before you move on}]
What does \kn{ಕನ್ನಡಿ} mean? (\textquotedblleft{}A mirror\textquotedblright{}.) Say it once more.
\end{tcolorbox}
\section[\texorpdfstring{\kn{ಟೀವಿ} (ṭīvi)}{ಟೀವಿ (ṭīvi)}]{\kn{ಟೀವಿ} (ṭīvi) — a television}

\label{lesson:KA-C137-tivi}



\begin{quote}
Before the new one: say the Kannada for a window, then the Kannada for a mirror.
\end{quote}

\subsection*{You'll want to know: \kn{ಟೀವಿ}}
\textbf{\kn{ಟೀವಿ}} — \emph{ṭīvi} — \textquotedblleft{}a television\textquotedblright{}.

\begin{grammarlens}[title={What you've built}]
One more word for this chapter.
\end{grammarlens}

\subsection*{Your turn}
\begin{itemize}
\raggedright
  \item \emph{Say it:} \emph{ṭīvi}
  \item \emph{Say it:} \emph{ṭīvi}, once more
  \item \emph{Say it:} say \emph{ṭīvi}
  \item \emph{Recall:} say the Kannada for a window, then the Kannada for a mirror, then say \emph{ṭīvi} again
\end{itemize}

\begin{tcolorbox}[breakable,colback=teal!4,colframe=teal!35!black,title={Before you move on}]
What does \kn{ಟೀವಿ} mean? (\textquotedblleft{}A television\textquotedblright{}.) Say it once more.
\end{tcolorbox}
\section[\texorpdfstring{\kn{ಕಂಪ್ಯೂಟರ್} (kampyūṭar)}{ಕಂಪ್ಯೂಟರ್ (kampyūṭar)}]{\kn{ಕಂಪ್ಯೂಟರ್} (kampyūṭar) — a computer}

\label{lesson:KA-C137-kampyutar}



\begin{quote}
Before the new one: say the Kannada for a mirror, then the Kannada for a television.
\end{quote}

\subsection*{You'll want to know: \kn{ಕಂಪ್ಯೂಟರ್}}
\textbf{\kn{ಕಂಪ್ಯೂಟರ್}} — \emph{kampyūṭar} — \textquotedblleft{}a computer\textquotedblright{}.

\begin{grammarlens}[title={What you've built}]
One more word for this chapter.
\end{grammarlens}

\subsection*{Your turn}
\begin{itemize}
\raggedright
  \item \emph{Say it:} \emph{kampyūṭar}
  \item \emph{Say it:} \emph{kampyūṭar}, once more
  \item \emph{Say it:} say \emph{kampyūṭar}
  \item \emph{Recall:} say the Kannada for a mirror, then the Kannada for a television, then say \emph{kampyūṭar} again
\end{itemize}

\begin{tcolorbox}[breakable,colback=teal!4,colframe=teal!35!black,title={Before you move on}]
What does \kn{ಕಂಪ್ಯೂಟರ್} mean? (\textquotedblleft{}A computer\textquotedblright{}.) Say it once more.
\end{tcolorbox}
